% ============================================================
% Title supplied by appendix wrapper.
\label{sec:advanced-operators-iv}
% ============================================================

These operators formalize auxiliary computational and algebraic structures
that build on the main spine's transport, recursion, and monodromy tools:
orthogonality enforcement (polar projection), shadow commutators
(non-commutation with observation), and interference kernels (averaged unitary
constraints).

% ------------------------------------------------------------
\subsection{Orthogonality Projectors (Nearest Unitary/Orthogonal Map)}
\label{subsec:orthogonality-projector}
% ------------------------------------------------------------

The orthogonality projector extracts the unitary (or orthogonal) factor from an operator via polar decomposition.

\begin{definition}[Orthogonality projector]
\label{def:orthogonality-projector}\lean{CoherenceCalculus.AppendixOperators.OrthogonalityProjector}
Let $A: \cH \to \cH$ be a bounded linear operator on a finite-dimensional inner-product space with polar decomposition $A = U |A|$, where $|A| = (A^* A)^{1/2}$ and $U$ is a partial isometry.
The \emph{orthogonality projector} is the map
\[
\mathcal{O}(A) := U.
\]
When $A$ is invertible, $\mathcal{O}(A)$ is unitary (or orthogonal over $\bbR$).
\end{definition}

\begin{proposition}[Procrustes optimality (Standard)]
\label{prop:procrustes}\lean{CoherenceCalculus.AppendixOperators.procrustes_optimality_statement}
For an invertible operator $A$ on a finite-dimensional real inner-product space, $\mathcal{O}(A)$ is the unique minimizer of
\[
\min_{Q \in O(n)} \|A - Q\|_F,
\]
where $O(n)$ denotes the orthogonal group and $\|\cdot\|_F$ is the Frobenius norm.
\end{proposition}

\begin{proof}[Proof (standard)]
This is the classical orthogonal Procrustes theorem.
Write $A = U |A|$ and use the fact that $\|A - Q\|_F^2 = \|A\|_F^2 + n - 2\langle Q, A \rangle_F$.
Maximizing $\langle Q, A \rangle_F = \Tr(Q^\top A)$ over orthogonal $Q$ gives $Q = U$.
\end{proof}

\begin{corollary}[Fixed points]
\label{cor:orthogonality-fixed}\lean{CoherenceCalculus.AppendixOperators.orthogonality_fixed_point}
If $A$ is already unitary (or orthogonal), then $\mathcal{O}(A) = A$.
In this case, $|A| = I$ and $A^{-1} = A^*$.
\end{corollary}

\begin{proof}
For unitary $A$, we have $A^* A = I$, so $|A| = I$ and $A = A \cdot I = \mathcal{O}(A) \cdot I$.
\end{proof}

\begin{remark}[Connection to aliasing operator]
\lean{CoherenceCalculus.AppendixOperators.PolarDecomp}
The orthogonality projector $\mathcal{O}$ is the same polar-decomposition extraction used for the aliasing operator $U_h$ (Section~\ref{subsec:aliasing}).
Both extract the ``direction'' (isometric part) from an operator, discarding the ``magnitude'' (positive part).
\end{remark}

\begin{example}[2$\times$2 orthogonality projection]
\lean{CoherenceCalculus.AppendixOperators.orthogonality_2x2_example}
\label{ex:orthogonality-2x2}
Let $A = \begin{psmallmatrix} 2 & 0 \\ 0 & 3 \end{psmallmatrix}$.
Then $|A| = A$ (since $A$ is positive) and $\mathcal{O}(A) = I$.
For $A = \begin{psmallmatrix} 0 & 2 \\ 3 & 0 \end{psmallmatrix}$, the SVD gives $\mathcal{O}(A) = \begin{psmallmatrix} 0 & 1 \\ 1 & 0 \end{psmallmatrix}$, the nearest orthogonal matrix.
\end{example}

\begin{quote}\small
\textbf{Intuition.}
The orthogonality projector ``enforces unitarity'' by replacing any operator with its closest unitary/orthogonal factor.
This is a constraint-satisfaction operator: given an approximate unitary, extract the exact unitary it approximates.
\end{quote}

% ------------------------------------------------------------
\subsection{Shadow Commutators (Noncommutation with Observation)}
\label{subsec:shadow-commutator}
% ------------------------------------------------------------

The shadow commutator measures the failure of an operator to commute with the horizon projection.

\begin{definition}[Shadow commutator]
\label{def:shadow-commutator}\lean{CoherenceCalculus.AppendixOperators.shadowCommutator}
For a bounded linear operator $A: \cH \to \cH$ and horizon projection $\CCPi{h}$, the \emph{shadow commutator} is
\[
\CCShadowComm{A}{h} := [A, \CCPi{h}] := A \CCPi{h} - \CCPi{h} A.
\]
\end{definition}

\begin{lemma}[Vanishing criterion]
\label{lem:shadow-vanish}\lean{CoherenceCalculus.AppendixOperators.shadowCommutator_vanish_iff}
$\CCShadowComm{A}{h} = 0$ if and only if $A$ preserves the horizon decomposition: $A$ maps $\cH_{\mathrm{obs}} = \CCPi{h} \cH$ to itself and $\cH_{\mathrm{sh}} = \CCPibar{h} \cH$ to itself.
\end{lemma}

\begin{proof}
$[A, \CCPi{h}] = 0$ means $A \CCPi{h} = \CCPi{h} A$.
For $x \in \cH_{\mathrm{obs}}$, we have $\CCPi{h} x = x$, so $A x = A \CCPi{h} x = \CCPi{h} A x \in \cH_{\mathrm{obs}}$.
For $y \in \cH_{\mathrm{sh}}$, we have $\CCPi{h} y = 0$, so $\CCPi{h} A y = A \CCPi{h} y = 0$, meaning $A y \in \cH_{\mathrm{sh}}$.
The converse follows by reversing the argument.
\end{proof}

\begin{lemma}[Decomposition into leakage components]
\label{lem:shadow-decomposition}\lean{CoherenceCalculus.AppendixOperators.shadowCommutator_decomposition_statement}
The shadow commutator decomposes as
\[
\CCShadowComm{A}{h} = \CCPibar{h} A \CCPi{h} - \CCPi{h} A \CCPibar{h}.
\]
In terms of block notation (observable/shadow), this is $A_{so} - A_{os}$ where $A_{so} = \CCPibar{h} A \CCPi{h}$ and $A_{os} = \CCPi{h} A \CCPibar{h}$.
\end{lemma}

\begin{proof}
Expand using $I = \CCPi{h} + \CCPibar{h}$:
\begin{align*}
A \CCPi{h} - \CCPi{h} A &= A \CCPi{h} - \CCPi{h} A (\CCPi{h} + \CCPibar{h}) \\
&= A \CCPi{h} - \CCPi{h} A \CCPi{h} - \CCPi{h} A \CCPibar{h} \\
&= (\CCPi{h} + \CCPibar{h}) A \CCPi{h} - \CCPi{h} A \CCPi{h} - \CCPi{h} A \CCPibar{h} \\
&= \CCPibar{h} A \CCPi{h} - \CCPi{h} A \CCPibar{h}. \qedhere
\end{align*}
\end{proof}

\begin{remark}[Connection to leakage operator]
\lean{CoherenceCalculus.AppendixOperators.shadowCommutator_decomposition_statement}
When $A = d$ is a nilpotent differential, the term $\CCPibar{h} d \CCPi{h} = \CCLeak{d}{h}$ is exactly the leakage operator (Definition~\ref{def:leakage-operator}).
Thus $\CCShadowComm{d}{h} = \CCLeak{d}{h} - d_{os}$, where $d_{os} = \CCPi{h} d \CCPibar{h}$ is the ``return from shadow'' component.
\end{remark}

\begin{remark}[Connection to coherence defect]
\lean{CoherenceCalculus.AppendixOperators.shadowCommutator}
For a linear operator $A$, the coherence defect (Definition~\ref{def:coherence-defect}) satisfies
\[
\CCTau{h}{A}{x} = \CCPi{h} A x - A \CCPi{h} x = -\CCPi{h} [A, \CCPi{h}] x = -\CCPi{h} \CCShadowComm{A}{h} x.
\]
Thus the observable part of the shadow commutator directly encodes the coherence defect.
\end{remark}

\begin{example}[Off-diagonal operator]
\lean{CoherenceCalculus.AppendixOperators.shadow_commutator_example}
\label{ex:shadow-commutator}
Let $\cH = \bbR^2$, $\CCPi{h} = e_1 e_1^*$ (projection onto first coordinate), and $A = \begin{psmallmatrix} 0 & 1 \\ 1 & 0 \end{psmallmatrix}$.
Then:
\[
\CCShadowComm{A}{h} = A \CCPi{h} - \CCPi{h} A = \begin{psmallmatrix} 0 & 0 \\ 1 & 0 \end{psmallmatrix} - \begin{psmallmatrix} 0 & 1 \\ 0 & 0 \end{psmallmatrix} = \begin{psmallmatrix} 0 & -1 \\ 1 & 0 \end{psmallmatrix}.
\]
The nonzero commutator shows $A$ does not preserve the horizon decomposition.
\end{example}

\begin{quote}\small
\textbf{Intuition.}
The shadow commutator $\CCShadowComm{A}{h} = [A, \CCPi{h}]$ measures ``horizon incompatibility.''
Information loss under observation is not erasure; it is exactly the failure of operators to commute with the horizon projection.
Operators with $\CCShadowComm{A}{h} = 0$ are ``horizon-compatible'' in the
sense that they respect the observable/shadow split.
\end{quote}

% ------------------------------------------------------------
\subsection{Interference Kernels (Averaged Unitary Constraints)}
\label{subsec:interference-kernel}
% ------------------------------------------------------------

An interference kernel averages a collection of unitary operators, with fixed points characterizing invariant states.

\begin{definition}[Interference kernel]
\label{def:interference-kernel}\lean{CoherenceCalculus.AppendixOperators.interferenceKernel4}
Let $\{U_1, U_2, \ldots, U_m\}$ be a finite collection of unitary operators on $\cH$.
The \emph{interference kernel} is the average
\[
\mathcal{K} := \frac{1}{m} \sum_{j=1}^{m} U_j.
\]
\end{definition}

\begin{definition}[Common fixed-point subspace]
\label{def:fixed-point-subspace}\lean{CoherenceCalculus.AppendixOperators.isCommonFixedPoint4}
The \emph{common fixed-point subspace} is
\[
\mathrm{Fix}(\{U_j\}) := \{ x \in \cH : U_j x = x \text{ for all } j = 1, \ldots, m \}.
\]
\end{definition}

\begin{lemma}[Fixed points are eigenvectors]
\label{lem:fixed-eigenvector}\lean{CoherenceCalculus.AppendixOperators.fixed_implies_kernel_eigenvector4}
If $x \in \mathrm{Fix}(\{U_j\})$, then $\mathcal{K} x = x$.
That is, $\mathrm{Fix}(\{U_j\}) \subseteq \Ker(\mathcal{K} - I)$.
\end{lemma}

\begin{proof}
If $U_j x = x$ for all $j$, then
\[
\mathcal{K} x = \frac{1}{m} \sum_{j=1}^{m} U_j x = \frac{1}{m} \sum_{j=1}^{m} x = x. \qedhere
\]
\end{proof}

\begin{lemma}[Spectral bound]
\label{lem:kernel-spectral}\lean{CoherenceCalculus.AppendixOperators.kernel_spectral_bound_statement}
The interference kernel satisfies $\|\mathcal{K}\| \le 1$.
Furthermore, $\mathcal{K} x = x$ with $\|x\| = 1$ implies $U_j x = x$ for all $j$ (when $\cH$ is finite-dimensional).
\end{lemma}

\begin{proof}
Since each $U_j$ is unitary, $\|U_j\| = 1$, so $\|\mathcal{K}\| \le \frac{1}{m} \sum_j \|U_j\| = 1$.

For the second statement: if $\mathcal{K} x = x$ with $\|x\| = 1$, then
\[
1 = \|x\|^2 = \langle x, \mathcal{K} x \rangle = \frac{1}{m} \sum_{j=1}^{m} \langle x, U_j x \rangle.
\]
Since $|\langle x, U_j x \rangle| \le \|x\| \|U_j x\| = 1$ for each $j$, and the average equals 1, each term must equal 1.
For unit vectors, $\langle x, U_j x \rangle = 1$ implies $U_j x = x$.
\end{proof}

\begin{corollary}[Eigenspace characterization]
\label{cor:kernel-eigenspace}\lean{CoherenceCalculus.AppendixOperators.kernel_eigenspace_forward4}
In finite dimensions, $\Ker(\mathcal{K} - I) = \mathrm{Fix}(\{U_j\})$.
\end{corollary}

\begin{proof}
Lemma~\ref{lem:fixed-eigenvector} gives one inclusion.
Lemma~\ref{lem:kernel-spectral} gives the reverse.
\end{proof}

\begin{remark}[Connection to representation projectors]
\lean{CoherenceCalculus.AppendixOperators.IsotypicProperties}
When $\{U_j\} = \{U(g) : g \in G\}$ for a finite group $G$, the interference kernel becomes
\[
\mathcal{K} = \frac{1}{|G|} \sum_{g \in G} U(g) = P_{\mathrm{triv}},
\]
the projector onto the trivial isotypic component (Section~\ref{subsec:representation-projector}).
The fixed-point subspace $\mathrm{Fix}(\{U(g)\})$ is precisely the $G$-invariant subspace.
\end{remark}

\begin{example}[Two reflections in $\bbR^2$]
\lean{CoherenceCalculus.AppendixOperators.interference_reflections_example}
\label{ex:interference-reflections}
Let $U_1 = I$ and $U_2 = \begin{psmallmatrix} 1 & 0 \\ 0 & -1 \end{psmallmatrix}$ (reflection across the $x$-axis).
Then:
\[
\mathcal{K} = \frac{1}{2}(U_1 + U_2) = \frac{1}{2}\begin{psmallmatrix} 2 & 0 \\ 0 & 0 \end{psmallmatrix} = \begin{psmallmatrix} 1 & 0 \\ 0 & 0 \end{psmallmatrix}.
\]
This is projection onto the $x$-axis, which is $\mathrm{Fix}(\{U_1, U_2\}) = \Span\{e_1\}$.
\end{example}

\begin{quote}\small
\textbf{Intuition.}
The interference kernel $\mathcal{K}$ averages unitary constraints.
Its eigenvalue-1 eigenspace is exactly the subspace satisfying all constraints simultaneously.
This provides a spectral characterization of ``satisfies all unitary constraints'' without iterative solving.
\end{quote}
